\begin{trapbox}
\textbf{易错点1：将“两条切线”误写成“两条有斜率的切线”。}

外点确实能引两条切线，但 $x_0=\pm a$ 时，一条是 $x=\pm a$，没有斜率。此时 (C009-1) 的 $k^2$ 系数为零，只保留另一条切线的有限斜率。

\textbf{易错点2：不检查韦达定理的使用条件。}

只有当 $a^2-x_0^2\ne0$ 且方程有两个实根时，才可以谈 $k_1+k_2$、$k_1k_2$。特别是过 $M(2,2)$ 引 $\frac{x^2}{4}+y^2=1$ 的切线，不能声称“斜率积为某个数”，因为一条斜率不存在。

\textbf{易错点3：误用判别式。}

公式
\[
\Delta_k=4a^2b^2\left(\frac{x_0^2}{a^2}+\frac{y_0^2}{b^2}-1\right)
\]
本身正确；但用“$\Delta_k>0$ 意味着两个实数斜率”的前提是 $a^2-x_0^2\ne0$。当它为零时，原式不是二次方程；要补看竖直线。

\textbf{易错点4：忽略椭圆上的特殊点。}

椭圆上非长轴顶点的切线只有一条，斜率为 $-\frac{b^2x_0}{a^2y_0}$；在 $P=(\pm a,0)$ 处，切线 $x=\pm a$ 的斜率不存在。“所有椭圆上点的切线斜率都属于 $\mathbb R$”是错误表述。

\textbf{易错点5：把固定椭圆的斜率全集和固定外点的斜率集合混为一谈。}

任意 $k\in\mathbb R$ 均可作为椭圆某条切线的斜率；但\textbf{固定}外点 $M$ 后，有限斜率只能取 (C009-1) 的实根，至多两个。

\textbf{易错点6：切线判定式没有考虑竖直直线。}

$m^2=a^2k^2+b^2$ 只用于 $y=kx+m$。若直线竖直，要直接判定 $x=\pm a$ 为椭圆切线，不应给竖直直线硬塞一个“无穷大斜率”。

\textbf{易错点7：蒙日圆证明漏掉一横一直的情况。}

写 $k_1k_2=-1$ 必须先确认两斜率都存在。若一个竖直，则另一个水平；相应的外点为 $(\pm a,\pm b)$，它们也在 $x_0^2+y_0^2=a^2+b^2$ 上。

\textbf{易错点8：把椭圆缩放成圆后，直接把角度也搬回来。}

坐标变换 $U=x/a,\ V=y/b$ 保持切点与切线的关系，因此能证明切线的条数；但当 $a\ne b$ 时，它一般\textbf{不保持垂直与角度}，不能直接用单位圆中的直角来证明椭圆的蒙日圆。

\textbf{考场自检四步：}先定位 $M$（内、上、外）；再看 $x_0=\pm a$ 否；然后列 $k$ 方程；最后才用判别式、韦达定理或夹角公式。
\end{trapbox}
